\section{Metodologia}\label{sec:metodo}

\subsection{Delineamento, unidades de análise e período}

Trata-se de pesquisa quantitativa, de natureza descritiva e explicativa, com delineamento longitudinal de estudo de casos múltiplos: dois municípios brasileiros (Manaus e Parintins) acompanhados ano a ano e um país de referência (Espanha), analisado por subsetor das administrações públicas. As unidades de análise são, portanto, apenas as três atribuídas no trabalho. Como cada município contribui com no máximo dez observações anuais, a análise privilegia a descrição detalhada das séries e usa as regressões como evidência complementar, com especificações simples, erros-padrão adequados a amostras pequenas e testes de sensibilidade explícitos.

O período é 2016--2025, os dez exercícios mais recentes com Declaração de Contas Anuais (DCA) completa no Siconfi na data da coleta (21 de outubro de 2026). As coberturas diferem entre variáveis e unidades, e essas diferências não foram preenchidas artificialmente: (i) Parintins não tem DCA de 2016 homologada no Siconfi --- a API do Tesouro não retorna registros nem para a DCA nem para o RREO do 6.º bimestre, e o extrato de entregas do ente registra para aquele ano apenas os relatórios de gestão fiscal ---, de modo que sua série começa em 2017; (ii) o PIB municipal do IBGE está disponível até 2023, o que limita as regressões que usam a renda local a 2016--2023 (15 observações: 8 de Manaus e 7 de Parintins); e (iii) para a Espanha, todas as séries do Eurostat estão disponíveis até 2025.

\subsection{Fontes e critérios de seleção dos dados}

Os dados municipais foram obtidos diretamente na API de dados abertos do Siconfi\footnote{SECRETARIA DO TESOURO NACIONAL. API de Dados Abertos do Siconfi -- Declaração de Contas Anuais (DCA), Anexos I-C, I-D e I-E. Disponível em: \url{https://apidatalake.tesouro.gov.br/ords/siconfi/tt/dca}. Acesso em: 21 out. 2026.}, que replica as Declarações de Contas Anuais enviadas pelos municípios. Foram usados os Anexos I-C (receitas orçamentárias por natureza), I-D (despesas por natureza) e I-E (despesas por função) da DCA, que consolidam o exercício completo, em vez do RREO, que pode ser retificado após o fechamento. A população, o PIB municipal e o IPCA vêm das tabelas 6579, 4714, 1378, 5938 e 1737 do Sidra/IBGE\footnote{INSTITUTO BRASILEIRO DE GEOGRAFIA E ESTATÍSTICA. Sistema IBGE de Recuperação Automática (Sidra): tabelas 1737, 6579, 4714, 1378 e 5938. Disponível em: \url{https://servicodados.ibge.gov.br/api/v3/agregados}. Acesso em: 21 out. 2026.}; os dados da Espanha vêm das tabelas \texttt{gov\_10a\_main}, \texttt{gov\_10dd\_edpt1}, \texttt{prc\_hicp\_aind}, \texttt{nama\_10\_gdp} e \texttt{demo\_pjan} do Eurostat\footnote{EUROSTAT. Tabelas gov\_10a\_main, gov\_10dd\_edpt1, prc\_hicp\_aind, nama\_10\_gdp e demo\_pjan. Disponível em: \url{https://ec.europa.eu/eurostat/api/dissemination/statistics/1.0/data/}. Acesso em: 21 out. 2026.}, compiladas segundo o Sistema Europeu de Contas (SEC 2010). Optou-se pelo Eurostat, e não pelo \textit{World Economic Outlook} do FMI utilizado no trabalho anterior, porque o Eurostat fornece os dados por subsetor (administração central, Comunidades Autônomas, entidades locais e seguridade social), indispensáveis para comparar níveis de governo, e porque suas séries são as notificadas no âmbito do procedimento de déficit excessivo. Todas as consultas, com URL, data e parâmetros, estão registradas nos arquivos brutos e no dicionário de variáveis do repositório.

\subsection{Variáveis e formas de mensuração}

O Quadro~\ref{tab:variaveis} apresenta as principais variáveis. O dicionário completo, com conceito, fonte, endereço eletrônico, período, cobertura, unidade, tratamento de ausentes e fórmula de cada variável, está no arquivo \texttt{04\_dicionario\_variaveis/dicionario\_variaveis.xlsx}.

\begin{table}[htbp]
\centering
\caption{Principais variáveis, definições e fontes}\label{tab:variaveis}
\scriptsize
\begin{tabular}{p{3.1cm}p{8.2cm}p{3.2cm}}
\toprule
Variável & Definição e mensuração & Fonte\\
\midrule
Receita total líquida & Receitas correntes + de capital, exceto intraorçamentárias, deduzidas as transferências ao FUNDEB e outras deduções (R\$) & Siconfi, DCA I-C\\
Receita tributária & Impostos, taxas e contribuições de melhoria (código 1.1), inclusive IRRF dos servidores & Siconfi, DCA I-C\\
Arrecadação própria corrente & Receita corrente bruta menos transferências correntes brutas & Siconfi, DCA I-C\\
Transferências líquidas & Transferências correntes (1.7) líquidas das deduções para o FUNDEB + transferências de capital (2.4) & Siconfi, DCA I-C\\
FPM, ICMS, FUNDEB, SUS & Rubricas identificadas pelo nome da conta nas três classificações vigentes (2016--2017, 2018--2021 e 2022--2025) & Siconfi, DCA I-C\\
Despesa empenhada & Despesas exceto intraorçamentárias, por natureza (pessoal, investimentos) e por função (saúde, educação) & Siconfi, DCA I-D e I-E\\
Dependência & Transferências líquidas / receita total líquida & Calculada\\
Autonomia tributária & Receita tributária / receita corrente líquida & Calculada\\
Resultado orçamentário & Receita total líquida menos despesa empenhada & Calculada\\
PIB per capita & PIB a preços correntes / população harmonizada & IBGE, tab. 5938\\
População harmonizada & Interpolação geométrica entre os Censos de 2010 e 2022 (2016--2021); Censo 2022; interpolação 2022--2024 (2023); estimativas pós-censo (2024--2025) & IBGE, tab. 1378, 4714 e 6579\\
Deflatores & IPCA, média anual do número-índice, base 2025 = 1; IPCH (HICP) da Espanha, média anual, base 2025 = 1 & IBGE tab. 1737; Eurostat\\
Espanha (por subsetor) & Receita total (TR), despesa total (TE), saldo (B.9), dívida (GD), impostos (D.2 + D.5 + D.91), transferências recebidas (D.7 + D.9), em \% do PIB e em euros de 2025 per capita & Eurostat\\
\bottomrule
\end{tabular}
\par\smallskip
\begin{minipage}{0.97\textwidth}\scriptsize Fonte: elaboração própria.\end{minipage}
\end{table}

\subsection{Tratamentos e transformações}\label{sec:tratamentos}

Cinco tratamentos merecem destaque, todos implementados em código e documentados no registro de alterações.

\textbf{Harmonização do plano de contas.} A classificação da receita mudou duas vezes no período: até 2017 vigorava a codificação herdada da Portaria Interministerial n.º 163/2001; entre 2018 e 2021, códigos de transição (por exemplo, 1.7.1.8.01.2.0 para o FPM); e, a partir de 2022, a estrutura definitiva (1.7.1.1.51.0.0 para o FPM). Como o mesmo conceito tem códigos diferentes em cada fase, as rubricas detalhadas foram identificadas pelo nome da conta, e não pelo código, com uma regra que soma apenas a conta de nível mais alto quando uma conta e suas subcontas atendem ao mesmo critério, evitando dupla contagem. As contas efetivamente somadas em cada ano estão listadas no arquivo de auditoria.

\textbf{Exclusão das operações intraorçamentárias.} Os agregados de despesa do Anexo I-D incluem as despesas intraorçamentárias (modalidade 91), sobretudo contribuições patronais ao regime próprio de previdência, que também aparecem como receita intraorçamentária. Para que receita e despesa sejam comparáveis, ambas foram tomadas exceto intraorçamentárias. Em Manaus, essa correção retira da despesa total entre R\$ 0,16 bilhão (2017) e R\$ 0,90 bilhão (2020) por ano, em valores nominais; Parintins não registra operações intraorçamentárias no período.

\textbf{Deduções do FUNDEB.} Os municípios contribuem compulsoriamente com 20\% de parte de suas transferências para o FUNDEB e recebem de volta valores redistribuídos dentro do estado\footnote{BRASIL. Lei n.º 14.113, de 25 de dezembro de 2020 (regulamenta o FUNDEB permanente). Disponível em: \url{https://www.planalto.gov.br/ccivil_03/_ato2019-2022/2020/lei/l14113.htm}.}. As receitas foram consideradas líquidas dessas deduções, e o FUNDEB recebido foi mantido entre as transferências, de modo que a dependência mede o que efetivamente financia o município.

\textbf{Deflação.} Os valores monetários brasileiros foram convertidos para reais de 2025 pelo fator $F_t = \bar I_{2025}/\bar I_t$, em que $\bar I_t$ é a média dos doze números-índices mensais do IPCA no ano $t$. Usa-se a média anual porque receitas e despesas são fluxos distribuídos ao longo do ano. Os valores espanhóis foram deflacionados da mesma forma pelo IPCH. O trabalho anterior usava fatores pontuais da Calculadora do Cidadão com base 2022 e atribuía fator diferente de 1 ao próprio ano-base, o que foi corrigido.

\textbf{População e valores per capita.} As estimativas anuais do IBGE para 2016--2021 foram projetadas a partir do Censo de 2010, e o Censo de 2022 revelou que elas superestimavam a população: em 2021, a estimativa oficial de Manaus excedia em \popDifMan\% o valor compatível com os dois censos, e a de Parintins, em \popDifPar\%. Usar as estimativas antigas e o Censo de 2022 na mesma série criaria uma quebra artificial nos valores per capita. Por isso, a série principal é uma população harmonizada, obtida por interpolação geométrica entre os censos (com datas de referência de 1.º de agosto para os censos e de 1.º de julho para as estimativas); para 2023, ano sem publicação oficial, usa-se a interpolação entre o Censo de 2022 e a estimativa de 2024. A série oficial é mantida na base e usada como teste de sensibilidade. Dados ausentes nunca foram substituídos por zero: permanecem como ausentes e as estatísticas usam apenas as observações disponíveis.

\subsection{Técnicas de análise}

A análise descritiva compara a evolução das receitas, das transferências e das despesas reais per capita de Manaus e Parintins entre si e com os subsetores espanhóis, e calcula, para cada município, a variação acumulada da despesa por real de variação das transferências entre o primeiro e o último ano com PIB disponível. Para testar as hipóteses H2 e H3, estima-se por mínimos quadrados ordinários (MQO), com os dados dos dois municípios empilhados (município-ano), o modelo
\begin{equation}\label{eq:fly}
G_{it} = \alpha + \delta P_i + \beta_1 T_{it} + \beta_2 Y_{it} + \varepsilon_{it},
\end{equation}
em que $G_{it}$ é a despesa empenhada real per capita do município $i$ no ano $t$, $T_{it}$ são as transferências líquidas reais per capita, $Y_{it}$ é o PIB real per capita (\textit{proxy} da renda local) e $P_i$ é uma variável indicadora de Parintins, que funciona como efeito fixo de município e absorve as diferenças permanentes de escala, de custo e de estrutura administrativa entre as duas cidades. Os coeficientes são propensões marginais a gastar: o teorema de equivalência prevê $\beta_1=\beta_2$, e o efeito \textit{flypaper} corresponde a $\beta_1>\beta_2$, com sinais esperados $\beta_1>0$ e $\beta_2>0$. A hipótese é testada pela estatística $t$ de $H_0\!:\beta_1-\beta_2=0$. A equação de esforço tributário é
\begin{equation}\label{eq:trib}
R_{it} = \alpha + \delta P_i + \gamma_1 T_{it} + \gamma_2 Y_{it} + u_{it},
\end{equation}
em que $R_{it}$ é a receita tributária real per capita. A hipótese de ilusão fiscal com substituição de esforço prevê $\gamma_1<0$; a ausência de efeito ($\gamma_1=0$) é compatível com transferências que não deslocam a arrecadação, como encontrado por \citeonline{litschig2013}.

Com 15 observações, os erros-padrão usuais e os assintóticos são pouco confiáveis. Por isso, usam-se erros robustos à heterocedasticidade do tipo HC3, os mais conservadores da família, e valores críticos da distribuição $t$ de Student com os graus de liberdade efetivos. Não se usam erros agrupados por município nem os de \citeonline{driscoll1998}, porque ambos exigem muitas unidades (ou muitos períodos) para serem consistentes, e aqui há apenas duas. A sensibilidade é avaliada por oito variações: (i) inclusão de tendência linear comum, que controla o crescimento simultâneo de receitas e despesas; (ii) separação do FPM, transferência incondicional típica, das demais transferências; (iii) especificação em logaritmos, cujos coeficientes são elasticidades; (iv) primeiras diferenças, que eliminam tendências e efeitos fixos; (v) valores per capita calculados com a população oficial do IBGE; (vi) e (vii) estimação separada para cada município; e (viii) exclusão do PIB, o que permite usar todos os anos até 2025 (19 observações). Os diagnósticos incluem fatores de inflação da variância (VIF), o teste de Breusch-Pagan para heterocedasticidade, o de Shapiro-Wilk para normalidade dos resíduos, o RESET de Ramsey para forma funcional, a autocorrelação de primeira ordem dos resíduos dentro de cada município (coeficiente AR(1) e estatística de Durbin-Watson) e a distância de Cook para observações influentes.

Para a Espanha, com dez anos por subsetor, estima-se uma regressão análoga, por MQO com erros robustos e efeito fixo de subsetor, da despesa real per capita das Comunidades Autônomas e das entidades locais sobre as transferências recebidas e o PIB per capita ($n=\esN$), lida apenas como indício descritivo da mesma regularidade.

\subsection{Associação, causalidade e comparabilidade}

Os coeficientes da Equação~\eqref{eq:fly} medem associações condicionais, e não efeitos causais. As transferências não são atribuídas aleatoriamente: o FPM depende da faixa populacional, a cota-parte do ICMS depende do valor adicionado e de critérios estaduais, e os repasses do SUS e do FUNDEB dependem de matrículas e de produção de serviços --- variáveis que também afetam o gasto. A variável indicadora de município remove as diferenças permanentes entre Manaus e Parintins, mas não as variações ao longo do tempo que afetam simultaneamente transferências e despesas, como o crescimento da população atendida ou mudanças nas regras de financiamento da saúde e da educação. Com apenas dois municípios, também não é possível controlar choques comuns por efeitos de ano, o que torna a tendência linear e as primeiras diferenças testes importantes. Além disso, parte das transferências é vinculada a gastos específicos, o que mecanicamente aproxima $\beta_1$ de 1 e torna o teste mais exigente de interpretar \cite{knight2002,megdal1987}. Desenhos quase-experimentais, como as descontinuidades do FPM \cite{brollo2013,litschig2013}, seriam necessários para uma interpretação causal; como discutido na Seção~\ref{sec:biblio}, esses desenhos têm suas próprias exigências \cite{baker2022,goodmanbacon2021}.

A comparação entre municípios e país também tem limites. As moedas, as escalas e os níveis de preços são diferentes, o que se contorna com proporções (da receita e do PIB) e com valores deflacionados por índices nacionais, sem conversão cambial. As contas brasileiras seguem a contabilidade orçamentária e patrimonial do MCASP\footnote{SECRETARIA DO TESOURO NACIONAL. \textit{Manual de Contabilidade Aplicada ao Setor Público (MCASP)}. 10.ª ed. Brasília: STN, 2023.}, de base mista, ao passo que as espanholas seguem o SEC 2010, de competência econômica. A organização fiscal também difere: as entidades locais espanholas têm competências mais restritas que os municípios brasileiros (saúde e educação são regionais), e o SEC registra como imposto regional a parcela cedida do IRPF e do IVA, que no Brasil seria classificada como transferência. Por isso, os municípios brasileiros são comparados com os dois subsetores espanhóis --- as Comunidades Autônomas, pela natureza das competências, e as entidades locais, pelo nível de governo ---, e as conclusões comparativas são qualitativas.

\subsection{Limitações metodológicas}

Além das já mencionadas, destacam-se: (i) a amostra é pequena (15 observações nas regressões com PIB), o que reduz o poder dos testes e torna os resultados sensíveis a observações individuais; (ii) o PIB per capita varia muito entre os dois municípios e pouco ao longo do tempo dentro de cada um, de modo que seu coeficiente é estimado com pouca precisão quando se controla o efeito fixo de município; (iii) o PIB municipal é uma \textit{proxy} imperfeita da renda dos residentes, sobretudo em Manaus, onde o valor adicionado industrial da Zona Franca não se converte integralmente em renda local, e foi deflacionado pelo IPCA por falta de deflator municipal; (iv) a população de 2023 é interpolada; (v) a qualidade da DCA depende do envio pelos municípios, e não foi possível auditar os lançamentos individuais; e (vi) com dez anos de dados, os testes para a Espanha também têm baixo poder estatístico. Esses limites não são contornados com imputações ou com a inclusão de outros municípios: a análise se restringe às unidades atribuídas, e as conclusões são formuladas na medida do que esses dados permitem afirmar.
